\documentclass[12pt]{article}
\usepackage[a4paper,margin=1in]{geometry}
\usepackage{amsmath,amssymb,amsfonts}
\usepackage{bm}
\usepackage{enumitem}
\usepackage{setspace}
\usepackage{titlesec}

\setstretch{1.15}
\setlength{\parskip}{0.5em}
\setlength{\parindent}{0pt}

\titleformat{\section}{\large\bfseries}{\thesection.}{0.5em}{}
\titleformat{\subsection}{\normalsize\bfseries}{\thesubsection}{0.5em}{}
\titleformat{\subsubsection}{\normalsize\bfseries}{\thesubsubsection}{0.5em}{}

\begin{document}

\begin{center}
{\LARGE \textbf{CSE400 -- Fundamentals of Probability in Computing}}\\[0.5em]
{\Large \textbf{Lecture 20: Linear Transformations and Expectations of Multiple Random Variables}}\\[0.5em]
\textbf{Source used:} Provided Lecture PDF only
\end{center}

\hrule
\vspace{1em}

\section{Outline of the Lecture}

The lecture covers the following topics exactly as listed in the provided material:

\begin{enumerate}
    \item Statistical Representation of Random Vectors
    \item Mean Vector, Correlation, and Covariance Matrices
    \item Linear Transformations
    \item Mapping properties under \( Y = AX + b \)
    \item Multivariate Gaussian Distributions
    \item From 1D/2D cases to N-dimensional Joint PDFs
    \item Special Cases and Properties
    \item Implications of zero correlation and matrix factorization
\end{enumerate}

\section{Random Vector Representation}

\subsection{Vector Form of Multiple Random Variables}

For random variables \( X_1, X_2, \ldots, X_N \), the lecture represents them in vector form as:

\[
X =
\begin{bmatrix}
X_1 \\
X_2 \\
\vdots \\
X_N
\end{bmatrix}
\]

\subsubsection*{Definition}

\begin{itemize}
    \item \(X\) is called a \textbf{random vector}.
    \item It provides a \textbf{compact representation of multiple random variables}.
\end{itemize}

\section{Mean Vector}

\subsection{Definition of Mean Vector}

For a random vector \(X\), the mean vector is defined as:

\[
\mu_X = E[X]
\]

\subsubsection*{Interpretation}

\begin{itemize}
    \item It represents the \textbf{average value} of the random vector.
    \item Each element corresponds to the mean of one random variable.
\end{itemize}

\subsection{Structure of the Mean Vector}

The lecture explicitly gives the structure as:

\[
\mu_X =
\begin{bmatrix}
E[X_1] \\
E[X_2] \\
\vdots \\
E[X_N]
\end{bmatrix}
\]

This shows that each coordinate of the mean vector is the expectation of the corresponding component random variable.

\section{Correlation Matrix}

\subsection{Definition}

The correlation matrix of the random vector \(X\) is defined as:

\[
R_{XX} = E[XX^T]
\]

The \((i,j)\)-th entry is:

\[
R_{XX}(i,j) = E[X_i X_j]
\]

where \(i\) and \(j < N\) as stated in the lecture slides.

\subsection{Matrix Representation}

The lecture writes the matrix explicitly as:

\[
R_{XX} =
\begin{bmatrix}
E[X_1X_1] & E[X_1X_2] & \cdots & E[X_1X_N] \\
E[X_2X_1] & E[X_2X_2] & \cdots & E[X_2X_N] \\
\vdots & \vdots & \ddots & \vdots \\
E[X_NX_1] & E[X_NX_2] & \cdots & E[X_NX_N]
\end{bmatrix}
\]

\subsubsection*{Property}

\begin{itemize}
    \item \(R_{XX}\) is a \textbf{symmetric matrix}.
\end{itemize}

\section{Covariance Matrix}

\subsection{Definition}

For a random vector \(X\), the covariance matrix is defined as:

\[
C_{XX} = E\left[(X - \mu_X)(X - \mu_X)^T\right]
\]

This is presented in the lecture as the ``relationship map'' of random vectors.

\subsection{Practical Form (2D Illustration)}

The lecture shows the 2-dimensional structure as:

\[
C_{XX} =
\begin{bmatrix}
\operatorname{Var}(X_1) & \operatorname{Cov}(X_1, X_2) \\
\operatorname{Cov}(X_2, X_1) & \operatorname{Var}(X_2)
\end{bmatrix}
\]

\subsubsection*{Symmetry Property}

The lecture explicitly notes:

\[
\operatorname{Cov}(X_i, X_j) = \operatorname{Cov}(X_j, X_i)
\]

Hence the covariance matrix is a \textbf{mirror of itself}, i.e., symmetric.

\section{Linear Transformation of Random Vectors}

\subsection{General Linear Model}

The lecture considers the transformation:

\[
Y = AX + b
\]

\subsubsection*{Interpretation of Terms}

\begin{itemize}
    \item \(A\): Transformation matrix
    \begin{itemize}
        \item Scales and rotates the data
    \end{itemize}
    \item \(b\): Constant vector
    \begin{itemize}
        \item Shifts / translates the center
    \end{itemize}
    \item \(X\): Input random vector
\end{itemize}

\subsection{Component-wise Form for N-dimensional Vectors}

The lecture writes the transformed components as:

\[
Y_1 = a_{11}X_1 + a_{12}X_2 + \cdots + a_{1N}X_N + b_1
\]

\[
Y_2 = a_{21}X_1 + a_{22}X_2 + \cdots + a_{2N}X_N + b_2
\]

\[
\vdots
\]

\[
Y_N = a_{N1}X_1 + a_{N2}X_2 + \cdots + a_{NN}X_N + b_N
\]

This explicitly shows how each output component is a linear combination of the input random variables plus a constant shift.

\section{Mean of the Transformed Random Vector}

\subsection{Goal}

For the linear transformation

\[
Y = AX + b
\]

the lecture defines the transformed mean as:

\[
\mu_Y = E[Y]
\]

\subsection{Step-by-Step Derivation (Exactly Following the Lecture)}

\subsubsection*{Step 1: Substitute the transformation into the expectation}

\[
\mu_Y = E[AX + b]
\]

\subsubsection*{Step 2: Apply linearity of expectation}

\[
\mu_Y = E[AX] + E[b]
\]

\subsubsection*{Step 3: Use the result stated in the lecture}

\[
\mu_Y = A\mu_X + b
\]

\subsection{Final Result}

\[
\boxed{\mu_Y = A\mu_X + b}
\]

This is the mean transformation rule given in the lecture.

\section{Correlation Matrix Transformation}

\subsection{Definition}

For the transformed random vector \(Y\), the lecture defines:

\[
R_{YY} = E[YY^T]
\]

Substituting \(Y = AX + b\):

\[
R_{YY} = E[(AX + b)(AX + b)^T]
\]

\subsection{Step-by-Step Derivation (Exactly Following the Lecture)}

\subsubsection*{Step 1: Start from the definition}

\[
R_{YY} = E[(AX + b)(AX + b)^T]
\]

\subsubsection*{Step 2: Expand the terms}

The lecture indicates ``Expansion of terms,'' followed by applying expectation.

\subsubsection*{Step 3: Apply expectation term-wise}

The lecture gives:

\[
R_{YY} = A E[XX^T] A^T + A E[X] b^T + b E[X^T] A^T + bb^T
\]

\subsubsection*{Step 4: Replace with standard notation}

Using

\[
R_{XX} = E[XX^T], \quad \mu_X = E[X]
\]

the lecture gives the final result:

\[
R_{YY} = A R_{XX} A^T + A\mu_X b^T + b\mu_X^T A^T + bb^T
\]

\subsection{Final Result}

\[
\boxed{R_{YY} = A R_{XX} A^T + A\mu_X b^T + b\mu_X^T A^T + bb^T}
\]

This is the correlation transformation rule presented in the lecture.

\section{Covariance Matrix Transformation}

\subsection{Key Observation}

The lecture explicitly states:

\[
Y - \mu_Y = (AX + b) - (A\mu_X + b) = A(X - \mu_X)
\]

\subsubsection*{Important Note}

\begin{itemize}
    \item The shift vector \(b\) \textbf{does not affect the deviation from the mean}.
    \item Therefore, the shift does not affect the spread.
\end{itemize}

\subsection{Step-by-Step Derivation (Exactly Following the Lecture)}

\subsubsection*{Step 1: Start with the covariance definition}

\[
C_{YY} = E[(Y - \mu_Y)(Y - \mu_Y)^T]
\]

\subsubsection*{Step 2: Substitute the deviation relation}

Since

\[
Y - \mu_Y = A(X - \mu_X)
\]

we get:

\[
C_{YY} = E\left[A(X - \mu_X)(X - \mu_X)^T A^T\right]
\]

\subsubsection*{Step 3: Use the simplified result stated in the lecture}

\[
C_{YY} = A C_{XX} A^T
\]

\subsection{Final Result}

\[
\boxed{C_{YY} = A C_{XX} A^T}
\]

\subsubsection*{Explicit Lecture Observation}

\begin{itemize}
    \item The shift \(b\) \textbf{does not affect the covariance (spread)}.
\end{itemize}

\section{Summary of Effects of Linear Transformation}

The lecture summarizes the three mapping properties under \(Y = AX + b\) as follows:

\subsection{Mean}

\[
\boxed{\mu_Y = A\mu_X + b}
\]

\subsection{Covariance}

\[
\boxed{C_{YY} = A C_{XX} A^T}
\]

\subsection{Correlation}

\[
\boxed{R_{YY} = A R_{XX} A^T + A\mu_X b^T + b\mu_X^T A^T + bb^T}
\]

These are the exact summary formulas shown in the lecture table.

\section{Worked Numerical Example: Linear Transformation}

\subsection{Given Data}

The lecture gives the transformation \(Y = AX\) with:

\[
\mu_X =
\begin{bmatrix}
2 \\
-1 \\
1
\end{bmatrix}
\]

\[
A =
\begin{bmatrix}
1 & 0 & -1 \\
0 & -1 & 1 \\
-1 & 1 & 0
\end{bmatrix}
\]

\[
R_{XX} =
\begin{bmatrix}
13 & 2 & 3 \\
2 & 10 & 3 \\
3 & 3 & 10
\end{bmatrix}
\]

The lecture asks to calculate:

\begin{enumerate}
    \item Mean vector of \(Y\)
    \item Correlation matrix of \(Y\)
    \item Covariance matrix of \(Y\)
\end{enumerate}

\subsection{Part (a): Mean Vector of \(Y\)}

\subsubsection*{Formula Used}

The lecture uses:

\[
\mu_Y = A\mu_X
\]

since the given example uses \(Y = AX\) (no shift vector shown).

\subsubsection*{Step-by-Step Calculation}

\[
\mu_Y =
\begin{bmatrix}
1 & 0 & -1 \\
0 & -1 & 1 \\
-1 & 1 & 0
\end{bmatrix}
\begin{bmatrix}
2 \\
-1 \\
1
\end{bmatrix}
\]

Compute each entry exactly as shown in the lecture:

\subsubsection*{First entry}

\[
(1)(2) + (0)(-1) + (-1)(1) = 2 + 0 - 1 = 1
\]

\subsubsection*{Second entry}

\[
(0)(2) + (-1)(-1) + (1)(1) = 0 + 1 + 1 = 2
\]

\subsubsection*{Third entry}

\[
(-1)(2) + (1)(-1) + (0)(1) = -2 - 1 + 0 = -3
\]

\subsubsection*{Final Mean Vector}

\[
\boxed{
\mu_Y =
\begin{bmatrix}
1 \\
2 \\
-3
\end{bmatrix}}
\]

This matches the mean vector used later in the covariance computation in the lecture.

\subsection{Part (b): Correlation Matrix of \(Y\)}

\subsubsection*{Formula Used}

The lecture states:

\[
R_{YY} = A R_{XX} A^T
\]

(Here the example uses \(Y=AX\), so only this term is used in the worked solution.)

\subsubsection*{Step 1: Compute \(A R_{XX}\)}

The lecture explicitly gives:

\[
A R_{XX} =
\begin{bmatrix}
10 & -1 & -7 \\
1 & -7 & 7 \\
-11 & 8 & 0
\end{bmatrix}
\]

This is shown as the result of the first multiplication step in the slides.

\subsubsection*{Step 2: Multiply by \(A^T\)}

The lecture then states the final result after multiplying by \(A^T\):

\[
\boxed{
R_{YY} =
\begin{bmatrix}
17 & -6 & -9 \\
-6 & 14 & -8 \\
-9 & -8 & 19
\end{bmatrix}}
\]

This is the correlation matrix shown in the worked example.

\subsection{Part (c): Covariance Matrix of \(Y\)}

\subsubsection*{Formula Used}

The lecture states:

\[
C_{YY} = R_{YY} - \mu_Y \mu_Y^T
\]

\subsubsection*{Step 1: Compute the outer product of the mean vector}

Using

\[
\mu_Y =
\begin{bmatrix}
1 \\
2 \\
-3
\end{bmatrix}
\]

the lecture gives:

\[
\mu_Y \mu_Y^T =
\begin{bmatrix}
1 \\
2 \\
-3
\end{bmatrix}
\begin{bmatrix}
1 & 2 & -3
\end{bmatrix}
=
\begin{bmatrix}
1 & 2 & -3 \\
2 & 4 & -6 \\
-3 & -6 & 9
\end{bmatrix}
\]

\subsubsection*{Step 2: Subtract from the correlation matrix}

The lecture writes:

\[
C_{YY} =
\begin{bmatrix}
17 & -6 & -9 \\
-6 & 14 & -8 \\
-9 & -8 & 19
\end{bmatrix}
-
\begin{bmatrix}
1 & 2 & -3 \\
2 & 4 & -6 \\
-3 & -6 & 9
\end{bmatrix}
\]

\subsubsection*{Step 3: Perform entry-wise subtraction}

\textbf{First row}

\[
17 - 1 = 16,\quad -6 - 2 = -8,\quad -9 - (-3) = -6
\]

\textbf{Second row}

\[
-6 - 2 = -8,\quad 14 - 4 = 10,\quad -8 - (-6) = -2
\]

\textbf{Third row}

\[
-9 - (-3) = -6,\quad -8 - (-6) = -2,\quad 19 - 9 = 10
\]

\subsubsection*{Final Result Shown in the Lecture}

The lecture slide explicitly displays:

\[
\boxed{
C_{YY} =
\begin{bmatrix}
16 & -8 & -8 \\
-8 & 10 & -2 \\
-8 & -2 & 10
\end{bmatrix}}
\]

This is the final covariance matrix shown in the provided PDF.

\section{Homework Assignment: Finding the Mean}

\subsection{Problem Statement}

The lecture provides:

\[
C_{XX} =
\begin{bmatrix}
5 & -1 & 2 \\
-1 & 5 & -2 \\
2 & -2 & 8
\end{bmatrix}
\]

\[
R_{XX} =
\begin{bmatrix}
6 & 1 & 1 \\
1 & 9 & -4 \\
1 & -4 & 9
\end{bmatrix}
\]

Task:

\begin{itemize}
    \item Find the mean vector \(\mu_X\).
\end{itemize}

\subsection{Hint Given in the Lecture}

The lecture uses the relationship:

\[
C_{XX} = R_{XX} - \mu_X \mu_X^T
\]

Rearranging:

\[
\mu_X \mu_X^T = R_{XX} - C_{XX}
\]

\subsubsection*{Note Given in the Lecture}

\begin{itemize}
    \item Identify the elements of \(\mu_X\) from the \textbf{diagonal} of the result.
\end{itemize}

\section{Gaussian Distributions: 1D and 2D}

\subsection{Univariate Gaussian (1D)}

The lecture introduces the case of a \textbf{single random variable \(X\)} under the heading:

\begin{itemize}
    \item ``1. Univariate Gaussian (1D)''
    \item ``For a single RV \(X\)''
\end{itemize}

(The explicit scalar formula is indicated in the slides, but only the heading and context are clearly readable in the provided excerpt.)

\subsection{Joint Gaussian (2D)}

For two random variables \(X\) and \(Y\), the lecture presents the \textbf{joint PDF} of the bivariate Gaussian distribution.

The slide explicitly identifies:

\begin{itemize}
    \item ``2. Joint Gaussian (2D)''
    \item ``For two RVs \(X\) and \(Y\), the joint PDF is:''
\end{itemize}

The displayed expression is the standard bivariate Gaussian form in terms of:

\begin{itemize}
    \item \(\mu_X, \mu_Y\)
    \item \(\sigma_X, \sigma_Y\)
    \item correlation coefficient \(\rho_{XY}\)
\end{itemize}

as shown on the lecture slide.

\section{Multivariate Gaussian Distribution (N-Dimensional)}

\subsection{Random Vector Form}

For an \(N\)-dimensional random vector, the lecture writes:

\[
X = [X_1, X_2, \ldots, X_N]^T
\]

\subsection{General Vector PDF}

The lecture presents the general multivariate Gaussian PDF as:

\[
f_X(x) =
\frac{1}{\sqrt{(2\pi)^N \det(C_{XX})}}
\exp\left(
-\frac{1}{2}(x-\mu_X)^T C_{XX}^{-1}(x-\mu_X)
\right)
\]

\subsubsection*{Structure Components Listed in the Lecture}

\begin{itemize}
    \item \(\mu_X\): Mean vector \((N \times 1)\)
    \item \(C_{XX}\): Covariance matrix \((N \times N)\)
\end{itemize}

\section{Example: 2D Joint Gaussian PDF}

The lecture includes a worked derivation for the 2D case by connecting the general multivariate form to the familiar bivariate Gaussian form.

\subsection{Step 1: Start from the General Multivariate Gaussian Form}

The lecture explicitly starts from:

\[
f_X(x) =
\frac{1}{\sqrt{(2\pi)^N \det(C_{XX})}}
\exp\left(
-\frac{1}{2}(x-\mu_X)^T C_{XX}^{-1}(x-\mu_X)
\right)
\]

and then specializes to the 2D case.

\subsection{Step 2: Compute the Determinant of the 2D Covariance Matrix}

The lecture explicitly states:

\[
\det(C_{XX}) = \sigma_1^2 \sigma_2^2 - (\rho \sigma_1 \sigma_2)^2
\]

Then simplifies it as:

\[
\det(C_{XX}) = \sigma_1^2 \sigma_2^2 (1 - \rho^2)
\]

This determinant is used in the normalization constant of the bivariate Gaussian density.

\subsection{Step 3: Inverse of the Covariance Matrix}

The lecture states that the inverse is obtained using the adjoint formula:

\[
C_{XX}^{-1} = \frac{1}{\det(C_{XX})}\operatorname{adj}(C_{XX})
\]

Then, after substitution and simplification, the lecture gives the standard 2D inverse form consistent with the bivariate Gaussian covariance structure.

\subsection{Step 4: Evaluate the Quadratic Form}

The lecture explicitly defines:

\[
Q(X) = (x-\mu_X)^T C_{XX}^{-1}(x-\mu_X)
\]

Then it states:

\begin{itemize}
    \item Substitute the 2D vector form into the quadratic expression.
    \item After algebraic simplification, the expression reduces to the standard bivariate Gaussian quadratic form.
\end{itemize}

\subsection{Conclusion of the Worked Derivation}

The purpose of the example is to show that the \textbf{general multivariate Gaussian formula} reduces to the \textbf{known bivariate Gaussian PDF} when:

\begin{itemize}
    \item \(N=2\),
    \item the covariance matrix is written in terms of \(\sigma_1^2, \sigma_2^2, \rho\).
\end{itemize}

\section{Example: Uncorrelated Gaussian Random Variables}

\subsection{Statement of the Example}

The lecture gives the following claim:

\begin{itemize}
    \item A vector of \(N\) jointly Gaussian random variables is \textbf{mutually uncorrelated}.
    \item Prove that the joint Gaussian PDF for a vector of uncorrelated random variables factorizes accordingly.
\end{itemize}

\subsection{Interpretation Given in the Lecture}

The lecture explicitly states:

\begin{enumerate}
    \item Since the variables are uncorrelated, the covariance matrix is \textbf{diagonal}.
    \item Therefore, the joint Gaussian PDF \textbf{factorizes into the product of individual Gaussian PDFs}.
\end{enumerate}

\subsection{Key Insight (Given in the Hint)}

For mutually uncorrelated Gaussian random variables:

\[
\operatorname{Cov}(X_i, X_j) = 0, \quad \forall i \neq j
\]

This is the exact condition highlighted in the lecture hint.

\subsection{Covariance Matrix Simplification}

Because all off-diagonal covariance terms are zero, the lecture shows that:

\[
C_{XX} =
\begin{bmatrix}
\sigma_1^2 & 0 & \cdots & 0 \\
0 & \sigma_2^2 & \cdots & 0 \\
\vdots & \vdots & \ddots & \vdots \\
0 & 0 & \cdots & \sigma_N^2
\end{bmatrix}
\]

Its inverse is therefore also diagonal:

\[
C_{XX}^{-1} =
\begin{bmatrix}
1/\sigma_1^2 & 0 & \cdots & 0 \\
0 & 1/\sigma_2^2 & \cdots & 0 \\
\vdots & \vdots & \ddots & \vdots \\
0 & 0 & \cdots & 1/\sigma_N^2
\end{bmatrix}
\]

This exact diagonal structure is shown in the lecture hint slides.

\subsection{Resulting Exponent}

The lecture explicitly notes:

\begin{itemize}
    \item Substitute the diagonal inverse back into the multivariate Gaussian density.
    \item The exponent becomes a sum of independent single-variable quadratic terms.
\end{itemize}

This is the reason the density factorizes into the product of 1D Gaussian PDFs.

\section{Homework Proof}

The lecture gives the following homework:

\[
\boxed{\operatorname{Var}[X+Y] = \operatorname{Var}[X] + \operatorname{Var}[Y] + 2\operatorname{Cov}(X,Y)}
\]

This is stated as a proof exercise only; no derivation is shown in the provided lecture slides. Therefore, no further steps are included here beyond the exact homework statement.

\section{Case Study: Smart City Traffic + Weather System}

The lecture concludes with an applied case study set in \textbf{Ahmedabad}, described as a system that predicts traffic congestion using transportation and climate variables.

\subsection{Step 1: Define \(N\) Random Variables}

The lecture defines a 5-dimensional random vector:

\[
X =
\begin{bmatrix}
X_1 \\
X_2 \\
X_3 \\
X_4 \\
X_5
\end{bmatrix}
\]

with the following meanings:

\begin{itemize}
    \item \(X_1\): Traffic flow (vehicles/min)
    \item \(X_2\): Average vehicle speed (km/h)
    \item \(X_3\): Rainfall intensity (mm/hr)
    \item \(X_4\): Temperature (\(^\circ\)C)
    \item \(X_5\): Air pollution index (AQI)
\end{itemize}

This is the system state random vector, with \(N=5\).

\subsection{Step 2: Mean Vector}

The lecture defines:

\[
\mu = E[X] =
\begin{bmatrix}
\mu_1 \\
\mu_2 \\
\mu_3 \\
\mu_4 \\
\mu_5
\end{bmatrix}
\]

and gives the following example average values:

\begin{itemize}
    \item Average traffic flow = 120 vehicles/min
    \item Average speed = 35 km/h
    \item Average rainfall = 2 mm/hr
    \item Average temperature = 32\(^\circ\)C
    \item Average AQI = 180
\end{itemize}

\subsection{Step 3: Covariance Matrix as a Dependency Map}

The lecture interprets the covariance matrix as the \textbf{dependency map of the city}.

It gives example sign interpretations:

\subsubsection*{Example 1}

\[
\operatorname{Cov}(X_1, X_3) > 0
\]

Interpretation:

\begin{itemize}
    \item More rain \(\Rightarrow\) more congestion (traffic increases)
\end{itemize}

\subsubsection*{Example 2}

\[
\operatorname{Cov}(X_2, X_3) < 0
\]

Interpretation:

\begin{itemize}
    \item More rain \(\Rightarrow\) lower speed
\end{itemize}

\subsubsection*{Example 3}

\[
\operatorname{Cov}(X_1, X_5) > 0
\]

Interpretation:

\begin{itemize}
    \item More traffic \(\Rightarrow\) higher pollution
\end{itemize}

\subsubsection*{Example 4}

\[
\operatorname{Cov}(X_4, X_5) < 0 \quad \text{(maybe)}
\]

Interpretation:

\begin{itemize}
    \item Higher temperature \(\Rightarrow\) pollution disperses
\end{itemize}

These interpretations are given directly in the lecture slides and are part of the intuition-building case study.

\subsection{Step 4: Joint Gaussian Assumption}

The lecture assumes:

\[
X \sim \mathcal{N}(\mu, \Sigma)
\]

\subsubsection*{Why Gaussian? (Exactly as Stated in the Lecture)}

\begin{enumerate}
    \item Aggregation of many small effects --- Central Limit Theorem
    \item Sensor noise is often Gaussian
    \item Easier analytics (closed-form solutions)
\end{enumerate}

\subsubsection*{Interpretation}

This means the entire system (traffic + climate) behaves in a structured probabilistic way defined by:

\begin{itemize}
    \item Mean (typical conditions)
    \item Covariance (interdependencies)
\end{itemize}

\subsection{Step 6: Linear Transformation to a Congestion Index}

The lecture defines a derived scalar metric:

\[
Y = 0.6X_1 - 0.4X_2 + 0.8X_3
\]

This is identified as a \textbf{linear transformation}.

The lecture writes:

\[
Y = a^T X
\]

where \(a\) is the coefficient vector corresponding to the linear combination.

\subsubsection*{Inference \#1 (Given in the Lecture)}

If

\[
X \sim \mathcal{N}(\mu,\Sigma)
\]

then:

\begin{enumerate}
    \item \(Y\) is also Gaussian
    \item Mean:
    \[
    E[Y] = a^T \mu
    \]
\end{enumerate}

The lecture indicates ``Variance:'' as the next item, but the explicit variance expression is not fully visible in the provided excerpt; therefore it is not expanded here beyond what is shown.

\subsection{Conceptual Summary from the Case Study}

The lecture concludes with the following conceptual mapping:

\begin{enumerate}
    \item Random vector = system state
    \item Covariance = relationships
    \item Gaussian = structured uncertainty
    \item Transformation = derived metrics
\end{enumerate}

\section{Final Exam-Oriented Summary}

The key examinable results from this lecture, exactly as presented, are:

\subsection{Random Vector}

\[
X =
\begin{bmatrix}
X_1\\
X_2\\
\vdots\\
X_N
\end{bmatrix}
\]

\subsection{Mean Vector}

\[
\mu_X = E[X]
\]

\subsection{Correlation Matrix}

\[
R_{XX} = E[XX^T]
\]

\subsection{Covariance Matrix}

\[
C_{XX} = E[(X-\mu_X)(X-\mu_X)^T]
\]

\subsection{Linear Transformation}

\[
Y = AX + b
\]

\subsection{Mean Under Linear Transformation}

\[
\boxed{\mu_Y = A\mu_X + b}
\]

\subsection{Correlation Under Linear Transformation}

\[
\boxed{R_{YY} = AR_{XX}A^T + A\mu_X b^T + b\mu_X^T A^T + bb^T}
\]

\subsection{Covariance Under Linear Transformation}

\[
\boxed{C_{YY} = AC_{XX}A^T}
\]

\subsection{Multivariate Gaussian PDF}

\[
\boxed{
f_X(x) =
\frac{1}{\sqrt{(2\pi)^N \det(C_{XX})}}
\exp\left(
-\frac{1}{2}(x-\mu_X)^T C_{XX}^{-1}(x-\mu_X)
\right)
}
\]

\subsection{Special Gaussian Property}

\begin{itemize}
    \item If jointly Gaussian variables are \textbf{mutually uncorrelated}, then:
    \begin{itemize}
        \item \(C_{XX}\) is diagonal
        \item The joint PDF factorizes into the product of individual Gaussian PDFs
    \end{itemize}
\end{itemize}

\end{document}